% ------------------------------------------------------------- Ethics
This work studies how an attacker can disguise interference with a UAV's
navigation as a benign disturbance. The method is dual-use, so we weighed its
benefits and harms before running any experiment, following the venue's ethics
guidelines~\cite{usenix25ethics}.

\paragraph{Stakeholders and harms.} The stakeholders are UAV operators, people
and property under flight paths, autopilot maintainers and regulators. The
potential harm is that an adversary uses the disguise taxonomy, or the
optimiser, to make a real spoofing attack harder to detect.

\paragraph{Why the benefit outweighs the risk.}
\begin{itemize}
\item \emph{The attack capability is not new.} Detector-agnostic stealthy GPS
  attacks against UAVs are already published~\cite{khazraei2024}, as are learned
  spoofers that exploit wind compensation~\cite{kim2024predictable}. Our
  contribution is the measurement of which disguise defeats which detector,
  together with a defence that does not depend on detection.
\item \emph{Every attack in this paper is paired with its mitigation in the
  same section:} re-anchoring-interval sizing, detector diversity across
  disguise classes, and the source-agnostic floor.
\item \emph{No experiment emitted radio-frequency energy.} All attacks act
  through simulator interfaces, and the released harness contains no
  transmit-capable code path.
\item \emph{No human subjects or personal data are involved.} The third-party
  corpora are public, are used under their licences, and are not
  redistributed.
\end{itemize}

\paragraph{Disclosure.} The pilot uses a generic Kalman filter rather than a
specific autopilot, so it discloses no vendor vulnerability. If the
software-in-the-loop transfer study in our runbook shows that a disguise family
evades PX4 EKF2's innovation gating, we will notify the PX4 and ArduPilot
security contacts at least 90 days before publication and include the
mitigation.
